\section{New nested-branch theorem and exact scope}
All statements in this new part use the designated input analytic quotient
and the original entire continuous class $\C$, energy $E$, raw mass two,
full bounded dual $D$, and actual decreasing-presentation class $\mathcal G$.
The full preceding 427 mathematical body is retained as data, with its
qualifications. The genuine nonalgebraization and coherent-descent premises
are not certified here. Neither an arbitrary Borel limit nor a classical
psh function is declared to belong to $\mathcal G$.

We prove a new equation sector consisting of arbitrary Radon distributions
along a nested nonskeletal bidisc branch, including distributions of
unbounded depth and unbounded finite-energy solutions. We also prove an
obstruction for an entire specified finite-fan local modification class and
all actual presentations using that class. Neither result resolves the
whole-mass input $2\delta_{x_{1/8}}$ in the entire original $\C/\mathcal G$.

\subsection{Applied foundations and what has actually been read}
The entire four actual five-field 427 notes, the distinct 431-LF/22233-byte
formal proof, all 2820 LF/137655 bytes of the standalone, the unique preamble
and abstract, current scope records, JSON and TXT ledgers, verdict, and
byte-unchanged residual were physically read. The four native programs were
read only as inert data. All fifteen retained immediate primary units were
freshly read in bounded complete outputs in this run. The initial truncated
inventories add no reading credit. The entire multi-megabyte historical
companion, deeper primary/book trees, geometry and descent have not been
freshly read or independently certified in their entirety.

Here are the exact external results applied, with their retained scopes.
CLD v2 Theorem 15.7.1 gives the calibrated tropical structure on a
topologically separated pure-dimensional analytic space with its actual
analytic boundary. Our input is pure dimension two and boundaryless; all
local current applications use ordinary open subsets. Definition 8.1.7
and Lemma 8.1.8 define the minimal continuous sheaf containing smooth
positive-Hessian functions and closed under local uniform limits.
Corollary 8.2.4 permits convex coordinatewise increasing composition.
Lemma 8.5.5 and Sections 8.7.1--2 give positive recursive products and
their local uniform continuity; 8.7.4 gives symmetry; 8.8.5 gives
multilinearity and invariance under smooth $\DD$-null changes. Proposition
8.5.6, iterated in 8.7.5, gives, for a topologically proper tropical map $j$
and a positive closed current $S$,
\[
 \DD u\wedge j_*S=j_*(\DD(j^*u)\wedge S).
\]
We apply this only to fixed closed analytic fibres over their fixed completed
residue fields. Propositions 7.2.8--13 give positive top Radon currents,
coefficient measures and the upgrade to continuous tests. Smooth Green and
Stokes are Proposition 5.11.1, with zero global boundary term here.
The whole-class strong-gradient, cubic, midpoint and fixed-test/full-point
moduli are reconstructed in the retained body, rather than inferred from a
verdict. Their conventional smooth-partition/graded/finite-flat/curve
foundations remain the stated audit boundary.

The precise fibre formula is CLD v2 Corollary 18.6.2: on a boundaryless,
topologically separated, equidimensional analytic space, for analytic
$f_1,\ldots,f_d$ and positive real $t,r_i$,
\[
 (\DD\log\max(t,|f_1|/r_1,\ldots,|f_d|/r_d))^d
 =\delta_{f^{-1}(\eta_{(tr_i)})}.
\]
The right hand side is its actual fibre integration current. Its proof uses
18.6.1 on the generically flat locus and computes the residual coefficient
as one. The literal first-case prose slips are not used as an assertion
that $\DD$ vanishes; the dimension argument concerns the top power.
Our coordinate translations, inversions and the monomial change
$(w_1,w_2)\mapsto(w_1/w_2,w_2)$ are actual isomorphisms, degree one.
We also use the $d=1$ formula first on the surface and then on its fixed
completed-residue-field curve. All radii are positive. No moving Gauss
section is asserted analytic and no additional extension degree occurs.

For the finite-fan obstruction, 14.5.2--3 gives primitive, saturated
degree-one calibration of an identity moment, and 5.10.1(c),(d)/5.10.2
gives calibrated integration and Radon reconstruction. The square in
dimension two has factor $2!$; the elementary convex area computation is
proved below. GR v2 unnumbered Section 1.2 supplies the CLD-to-classical
bridge; Definition 4.1 and Proposition 4.2 retain USC and every curve after
every complete analytic extension. Lemma 6.3 and Theorem 6.4 give finite
full-dimensional affinoid Shilov boundaries and their maximum principle.
They are used on continuous local potentials of actual presentations.
GR 6.5 is not a recognition input, and CLD 8.6.7 is not applied. The local
full-point modulus uses the finite-flat representative from CLD11.6.3,
11.6.5(c), and the graded basis criterion13.8.1--5, with the retained
conventional Berkovich/Temkin/Ducros foundations. None of these readings
constitutes a correctness grant or a complete recursive source audit.

\subsection{Actual weights and the newly computed mixed products}
Use the normalized baseline $q_0$ and the injective chart
$U=\{|r_i|<1/4\}$ from the retained construction. On this chart let
$w_i=z_i-1$, $A_i=\max(1,|z_i|)=\max(1,|w_i|)$, and set
\[
 L_c=\log\max(e^{-c},|w_1|,|w_2|,|w_1w_2|),\qquad
 v_c=L_c-\log(A_1A_2),\quad q_c=q_0+v_c\quad(c>0).
\]
Extend $v_c$ by zero outside $U$. Let $g=\sigma(0)$, let $x_c$ be the
joint Gauss point of the bidisc $|w_i|\le e^{-c}$, and put $x_0=g$.
These points depend continuously on $c\ge0$. Define the extended depth
function on $U$ by
\[
 \gamma(y)=\max\{0,-\log\max(|w_1(y)|,|w_2(y)|)\},
\]
and by zero outside $U$. It is continuous into $[0,+\infty]$, is infinite
only at the actual rigid point $P=(1,1)$ in this chart, and satisfies
$\gamma(x_c)=c$. Indeed, whenever one $|w_i|\ge1$ the correction is zero;
otherwise $v_c=-\min(c,\gamma)$. If some $|z_i|\ne1$, one $|w_i|\ge1$.
Thus the correction vanishes on the ordinary open collar
$1/8<\max|r_i|<1/4$. This proves genuine gluing on the quotient; $z_i-1$
is used only on the injective chart, so the deck units $2,3$ are not
silently discarded.

The local potential is $L_c-\frac12\log|z_1z_2|-1/4$. The positive
constant $e^{-c}$ permits omission of every vanishing polynomial slot on
an ordinary neighborhood. The remaining slots are logarithms of units;
smooth convex maxima converge locally uniformly and have positive Hessian.
The subtracted logarithms are smooth $\DD$-null terms. Consequently
$q_c\in\C$ by the entire minimal-sheaf definition, including at zeros.
The retained fixed-coordinate computation gives
\[
 \MA(q_0)=2\delta_g,\qquad \MA(q_c)=\delta_g+\delta_{x_c},
 \qquad B(q_0,q_c)=2\delta_g.
\]
For completeness the two vertices of $L_c$ have tropical radii
$(-c,-c)$ and $(0,0)$. The first is the translation fibre with coefficient
one; at the second subtract $\log|w_1w_2|$ and use the inverse-coordinate
fibre, also coefficient one. Elsewhere the square is one-dimensional or
affine and zero. The baseline square is the sum of two degree-one mixed
coordinate fibres. Restricting $L_c$ to either fixed $z_i$ unit Gauss fibre
gives $\log\max(1,|z_j|)$, proving the mixed baseline formula on actual
analytic tests. This reconstructs the previously supplied products.

\begin{lemma}\label{branch:mixed}
The new mixed formula, for every $0<c\le d$, is
\[
 B(q_c,q_d)=\delta_g+\delta_{x_c}.
\]
It is an equality of actual recursive currents and Radon measures on $X$.
\end{lemma}
\begin{proof}
For $c=d$ it is the square just computed. Suppose $c<d$.
Off the chart the weights agree with $q_0$ and its top measure is zero.
At polynomial zeros omit vanishing slots; both local potentials then
depend on one coordinate, so their mixed product is zero.
On the remaining open set use $b_i=\log|w_i|$.
If one $b_i\ge0$, the constant slots are strictly smaller and $L_c,L_d$
coincide locally with $\max(b_1,b_2,b_1+b_2)$; their product is the
previous square, supported only at the inverse-coordinate Gauss point $g$.
If both $b_i<0$, write $m=\max(b_1,b_2)$. When $m<-c$, $L_c$ is constant.
When $m>-c$, both equal $m$ locally; their square has only one unit ratio
and vanishes. When $m=-c$ and $b_1\ne b_2$, $L_d$ is a single unit
logarithm locally and its Hessian vanishes. Thus the only additional
possibility is $b_1=b_2=-c$.

Near this last set, $L_d=\log\max(|w_1|,|w_2|)$, hence
\[
 \DD L_d=\DD\log\max(1,|w_1/w_2|)
           =\delta_{(w_1/w_2)^{-1}(\eta_1)}.
\]
This is the $d=1$ formula on an ordinary open surface. On its closed fixed
fibre over $\mathscr H(\eta_1)$, the coefficient $w_1/w_2$ has constant
norm one, and $L_c$ restricts to $\log\max(e^{-c},|w_2|)$.
Fixed-current restriction and the curve $d=1$ formula give one point.
It is exactly $x_c$: the exponent map $(i,j)\mapsto(i,i+j)$ is a bijection
of $\mathbb Z^2$, so the iterated Gauss norm in $(w_1/w_2,w_2)$ equals
the joint Gauss norm of radius $(e^{-c},e^{-c})$ in $(w_1,w_2)$.
The determinant is one and the analytic coordinate change is an
isomorphism. This also excludes every non-Gauss point at the same radii.
At $g$ the common inverse-coordinate calculation gives coefficient one.
All maps, fibres and restrictions retain the source boundary hypotheses;
positivity and7.2.8--13 upgrade the identities to every continuous test.
\end{proof}

\begin{theorem}\label{branch:finite}
Let $0<c_1<\cdots<c_N$, $m_k>0$, and $M=\sum_km_k\le1$.
Put $\alpha_k=\sqrt{\sum_{j=k}^Nm_j}$, $\alpha_{N+1}=0$,
$\beta_k=\alpha_k-\alpha_{k+1}$, and
\[
 q=q_0+\sum_{k=1}^N\beta_kv_{c_k},\qquad
 d(c)=\sum_k\beta_k\min(c,c_k),\qquad S=\sum_km_kd(c_k).
\]
Then $q\in\C$ is normalized and
\[
 \MA(q)=(2-M)\delta_g+\sum_km_k\delta_{x_{c_k}},\qquad
 E(q)=-\tfrac13-\tfrac S3,\quad
 D(\MA(q))=\tfrac16+\tfrac{2S}3.
\]
The full dual is attained and $q+2^{-n}\downarrow q$ is an actual
finite-lower-energy presentation. Zero masses may be omitted; $M=0$
gives $q_0$. This is already larger than the one-depth 427 sector.
\end{theorem}
\begin{proof}
The coefficients are nonnegative and sum to $\alpha_1=\sqrt M\le1$;
thus $q=(1-\alpha_1)q_0+\sum\beta_kq_{c_k}$ is an actual positive
convex combination. It is at most $q_0$ and equals it at its maximum,
so its supremum is zero and its range is in $[-1/4-d(c_N),0]$.
By Lemma\ref{branch:mixed} and the baseline mixed formula, the coefficient
at $x_{c_k}$ is
\[
 \beta_k^2+2\beta_k\sum_{j>k}\beta_j
 =\alpha_k^2-\alpha_{k+1}^2=m_k.
\]
All other mass is at $g$, and the coefficient there is $2-\alpha_1^2$,
either by expansion or the common raw mass two. This is actual product
multilinearity, not an assumed equality of tropical projections.
For $v=q-q_0$, the same calculation at parameter $h\in[0,1]$ gives
$\MA(q_0+hv)=(2-h^2M)\delta_g+h^2\sum m_k\delta_{x_{c_k}}$.
Here $v(g)=0$ and $v(x_{c_k})=-d(c_k)$. The actual first variation is
$\frac{d}{dh}E(q_0+hv)=-h^2S$. The baseline energy is $-1/3$,
so integration gives the energy formula. The target moment is $-1/2-S$.
Entire-class tangent concavity gives the full-dual formula and attainment.
The stated sequence is continuous, decreasing and bounded in energy.
\end{proof}

\subsection{Every Radon depth distribution and its exact energy criterion}
Let $\zeta$ be any finite positive Radon measure on $(0,\infty)$, with
$M=\zeta((0,\infty))\le1$. Define
\[
 T(t)=\zeta((t,\infty)),\quad \alpha(t)=\sqrt{T(t)},\quad
 d_C(c)=\int_0^{\min(c,C)}\alpha(t)\,dt,\quad
 S_C=\int_0^C T(t)^{3/2}\,dt,
\]
and $S=\lim_CS_C\in[0,\infty]$, $d(c)=\int_0^c\alpha(t)dt$.
Endpoints at atoms do not affect these Lebesgue integrals.
Define the genuine Radon target on $X$ by
\[
 \mu_\zeta=(2-M)\delta_g+(c\mapsto x_c)_*\zeta.
\]
The pushforward is a finite regular Borel measure: restrict to successive
compact depth intervals, use continuous pushforward there, and approximate
in total variation by these restrictions. It has raw mass two.

\begin{theorem}\label{branch:radon}
Exactly for these targets one has the equivalence
\[
 D(\mu_\zeta)<\infty\quad\Longleftrightarrow\quad
 S=\int_0^\infty\zeta((t,\infty))^{3/2}dt<\infty.
\]
If $S<\infty$, the actual weight
\[
 q_\zeta(y)=q_0(y)-d(\gamma(y))
\]
has the actual continuous decreasing presentation
$q_n=q_0-d_n(\gamma)$ in the original $\mathcal G$, and
\[
 \MA(q_\zeta)=\mu_\zeta,\quad E(q_\zeta)=-\tfrac13-\tfrac S3,
 \quad\int q_\zeta\,d\mu_\zeta=-\tfrac12-S,
 \quad D(\mu_\zeta)=\tfrac16+\tfrac{2S}3.
\]
It is unique modulo constants among all actual original $\mathcal G$
solutions. It is bounded and continuous exactly when
$\int_0^\infty\sqrt{T(t)}dt<\infty$; otherwise it is continuous into
$\mathbb R\cup\{-\infty\}$, finite and locally continuous off $P$,
with value $-\infty$ at $P$. No mass at $P$ is introduced.
\end{theorem}
\begin{proof}
For fixed finite $C>0$, push $\zeta$ forward by $c\mapsto\min(c,C)$
to a measure $\zeta_C$ on $(0,C]$ of the same mass. Approximate it by
finite atomic measures of mass $M$, moving each small partition interval's
mass to its positive right endpoint. Continuous tests on $[0,C]$ converge
uniformly by their moduli of continuity. At each non-atomic threshold the
tails converge; a finite measure has only countably many atoms. Dominated
convergence therefore gives $L^1(0,C)$ convergence of their square-root
tails to $\alpha$. For an atomic approximation, its finite solution in
Theorem\ref{branch:finite} equals $q_0-\int_0^{\min(\gamma,C)}\alpha_j(t)dt$:
differentiate the finite sum of $\beta_k\min(c,c_k)$ between thresholds,
and compare its value at zero. Its uniform difference from
$q_C=q_0-d_C(\gamma)$ is at most $\int_0^C|\alpha_j-\alpha|dt$.
Thus $q_C\in\C$ by actual uniform closure, with the genuine ordinary
collar gluing already verified. Recursive product continuity and the
continuous Gauss branch give, on every continuous analytic test,
\[
 \MA(q_C)=(2-M)\delta_g+(c\mapsto x_{\min(c,C)})_*\zeta.
\]
These are genuine recursive products, independently of the weak target
calculation. In particular no function is recognized solely by its image
on a skeleton or its point tuple.

Tonelli gives the key exact identity
\[
 \int d_C(c)\,d\zeta(c)
 =\int_0^C\alpha(t)\zeta([t,\infty))dt=S_C.
\]
The equality of the two tails holds for almost every $t$. Passing the
finite energies and target moments to the uniform limit, or using actual
first variation along $q_0+h(q_C-q_0)$, now yields
\[
 E(q_C)=-1/3-S_C/3,\quad \int q_C\MA(q_C)=-1/2-S_C,
 \quad D(\MA(q_C))=1/6+2S_C/3.
\]
Moreover $q_C(x_c)=-1/4-d_C(c)$ for every finite $c$, so
$\int q_C\mu_\zeta=-1/2-S_C$ as well.

The $q_n$ are actual continuous nonpositive normalized weights and decrease
pointwise to the displayed $q_\zeta$. If $S<\infty$, their energies have a
finite lower bound. This is an actual original $\mathcal G$ presentation,
not an assignment of energy to a Borel relaxation. Its canonical energy
is the stated limit. The bounded MA measures converge to $\mu_\zeta$:
for any continuous $h$ the discrepancy is bounded by
$2\|h\|_\infty\zeta((n,\infty))$. Hence the original generated operator
has the asserted value. Tonelli with $C\to\infty$ proves the finite target
moment; $\zeta$ has no point at infinite depth.

For every fixed $a\in\C$, tangent concavity at $q_n$ gives
$E(a)-\int a\MA(q_n)\le1/6+2S_n/3$. Pass to the weak limit to bound the
full $D(\mu_\zeta)$ above. The actual bounded competitors $q_n$ give the
reverse bound because $E(q_n)-\int q_n\mu_\zeta=1/6+2S_n/3$.
This also proves $D(\mu_\zeta)=+\infty$ if $S=+\infty$, without constructing
an ineligible operator. Boundedness is equivalent to $d(\infty)<\infty$.
If finite, the convergence is uniform on all $X$; otherwise $d(\gamma)$
is finite and continuous off $P$, and tends to infinity on ordinary
neighborhoods shrinking to $P$. Thus the extended continuity is genuine.
Generated uniqueness is proved next, without assuming general recognition.
\end{proof}

\subsection{Canonical nonordered representation in this larger sector}
\begin{theorem}\label{branch:recognition}
For $S<\infty$, every normalized full-dual maximizing family for
$\mu_\zeta$ has the full-point limits of the actual weight $q_\zeta$,
canonical energy $E(q_\zeta)$, genuine weak MA limit $\mu_\zeta$, and
canonical ordinary upper relaxation $q_\zeta$. It converges in
$L^1(\mu_\zeta)$ to $q_\zeta$. Every selected fast sequence has $J=K=0$.
No positivity is assigned to its raw Borel limit at remaining points.
\end{theorem}
\begin{proof}
The $q_n$ maximize the full dual, with exact deficits
$2(S-S_n)/3$. Full-dual coercivity puts them and every late maximizing
family in a common normalized energy sublevel. The retained whole-class
midpoint formula gives
\[
 I(a_i,q_n)\le4\bigl(\varepsilon_i+2(S-S_n)/3\bigr).
\]
The actual energy, fixed-test, full-point and weak-test moduli, followed
first by large $n$ and then by late $i$, identify the claimed real energy,
each full-point value and the genuine weak measure limit. Every full point
is different from the rigid $P$, so $q_n(x)=q_\zeta(x)$ for all sufficiently
large $n$. For fixed continuous $|a_i-a_j|$, pass the uniform third-slot
bound against the energy-bounded $\MA(q_n)$ to its weak limit. This gives
$\int|a_i-a_j|\mu_\zeta\le T_L(I(a_i,a_j))$ for one common sublevel $L$.
Completeness gives a canonical $L^1(\mu_\zeta)$ limit for every net;
product-index comparison identifies it with the limit of $q_n$.
That limit is $q_\zeta$, by monotone convergence and its finite target
moment. No countable cofinality of a maximizing net is used.

The all-affinoid Shilov formula for the ordinary upper relaxation uses
only full-point values. On $X\setminus\{P\}$ the already constructed
$q_\zeta$ is continuous, so the Shilov maximum and shrinking ordinary
neighborhoods identify the relaxation with it. If $q_\zeta(P)$ is finite,
it is continuous at $P$ too and the same argument applies. Otherwise
$q_\zeta<-A$ on some ordinary neighborhood of $P$ for every finite $A$;
all its full-point values there obey that bound. The Shilov formula
forces the relaxation at $P$ to be $-\infty$. This proves actual
representation of the canonical relaxation by the previously constructed
$\mathcal G$ weight, including the unbounded case.
The selected fast limit equals $q_\zeta$ at every full point, hence at
$g$ and every $x_c$ in the target. Its target moment is the stated moment;
the exact loss identity yields $J=0$, and the weak measure equality yields
$K=0$. Other raw Borel points receive no declared MA operator.
\end{proof}

Here is the generated uniqueness argument, including its actual antecedent.
For any actual $b_n\downarrow v$ with finite lower energy, compare two
presentations by compact strict domination plus arbitrary constants.
Energy monotonicity and $E(a+\epsilon)=E(a)+2\epsilon$ make the limiting
energy canonical. Apply this also to each finite convex combination and
use finite unisolvent cubic interpolation to get canonical coefficient
limits. If $\lambda=\MA(v)$, actual first variation of
$E((1-t)b_n+tw)$ and $E(tb_n)$ gives the canonical weak limit against each
fixed $w\in\C$. The uniformly dense real span and common positive mass
extend it to all continuous tests. Let $e=\lim E(b_n)$ and
$p=\lim\int b_n\MA(b_n)$. Passing tangent inequalities gives
$D(\lambda)\le e-p$. After a common upper shift, the full dual inequality
and sequential monotone convergence show $v$ is $\lambda$-integrable.
The ordered cubic cocycle gives
\[
 0\le\int(b_n-b_l)\MA(b_l)\le3(E(b_n)-E(b_l))\quad(l\ge n).
\]
Pass $l$ against fixed continuous $b_n$, then $n$ by integrability:
$p=\int v\lambda$, and $D(\lambda)=e-\int v\lambda$.
Thus the normalized presentation maximizes this full dual. Its decreasing
suprema have finite limit $s=\sup v$ by nested compact superlevel sets
and $E(b_n)\le2\sup b_n$.

If $\lambda=\mu_\zeta$, Theorem\ref{branch:recognition} gives
$v-s=q_\zeta$ at all full points. These are dense by the nonempty finite
full-dimensional Shilov boundary in every small pure affinoid neighborhood.
On $X\setminus\{P\}$, USC of the already decreasing $v-s$ and continuity
of $q_\zeta$ prevent values below it. The finite Shilov maximum for each
continuous $\rho+b_n-\sup b_n$, followed by $n\to\infty$ and shrinking
neighborhoods, prevents values above it. At $P$ the same argument works
if $q_\zeta$ is continuous finite; if it is $-\infty$, its arbitrarily
negative bounds near $P$ force $v(P)=-\infty$. Hence equality holds at
every actual point. This is uniqueness modulo constants in all original
actual $\mathcal G$, without manufacturing a presentation for an unrelated
measure.

\begin{corollary}\label{branch:unboundedexample}
The explicit nonskeletal raw-mass-two Radon target
\[
 \mu=\delta_g+\int_1^\infty\delta_{x_c}\,c^{-2}dc
\]
has an unbounded actual finite-energy solution with
$E=-4/3$ and full $D=13/6$. Its correction is $-\gamma$ for
$0\le\gamma\le1$ and $-(2\sqrt\gamma-1)$ for $\gamma\ge1$, with value
$-\infty$ at $P$. For integers $n\ge1$, its actual continuous decreasing
presentation has energy $-4/3+2/(3\sqrt n)$.
\end{corollary}
\begin{proof}
Its depth tail is one for $0\le t<1$ and $t^{-1}$ for $t\ge1$.
Thus $S=1+\int_1^\infty t^{-3/2}dt=3$, whereas
$\int\sqrt{T}=\infty$. Integrate the square-root tail to get the
displayed correction and substitute $S_n=3-2/\sqrt n$ in the formulas.
All operator and uniqueness conclusions are the preceding theorem's.
\end{proof}

\begin{proposition}\label{branch:diracobjective}
For the whole Dirac target $2\delta_{x_a}$, $a>0$, the supremum of its
objective over the entire newly constructed branch-solution family
($M\le1$, $S<\infty$), or over all its bounded cutoff presentations, is
exactly $1/6+5a/3$. It is attained by $\zeta=\delta_a$, hence by $q_a$.
At $a=1/8$ this is $3/8$. This computes a family-restricted supremum,
not the original full dual of the whole Dirac.
\end{proposition}
\begin{proof}
At $x_a$ the constructed weight is $-1/4-d(a)$, so its objective is
\[
 1/6+2\int_0^a\alpha(t)dt-\tfrac13\int_0^\infty\alpha(t)^3dt.
\]
Here $0\le\alpha\le1$. On $[0,a]$ the function
$2\alpha-\alpha^3/3$ is strictly increasing, with maximum $5/3$;
outside this interval the integrand is nonpositive. Thus the expression
is at most $1/6+5a/3$, with equality for the tail of $\delta_a$.
Every cutoff simply replaces $\alpha$ by its restriction to $[0,n]$,
so the same bound holds for actual bounded competitors. No scalar
expression here is substituted for an operator on a Borel function.
\end{proof}

\subsection{A complete obstruction for the whole translated finite-fan class}
This obstruction applies to a specified local class, not all $\C$.
Let $H_0(b)=\max(0,b_1)+\max(0,b_2)$. Consider every finite real affine
maximum $H$ on $\mathbb R^2$ which is convex, contains a constant lower
slot, and agrees with $H_0$ whenever $b_1\ge0$ or $b_2\ge0$.
Use the potential $H(\log|w_1|,\log|w_2|)-\frac12\log|z_1z_2|-1/4$
on $U$ and the baseline outside, omitting vanishing slots as above.
Denote its global weight by $q_H$. No hypothesis of integral slopes is
made; these are real affine maxima of unit logarithms.

\begin{lemma}\label{branch:area}
For a finite convex affine maximum on a two-dimensional identity torus
moment, its actual Hessian square at the joint Gauss point with radii
$b_*$ has coefficient $2\operatorname{Area}(\partial H(b_*))$.
There is no additional face or analytic multiplicity for these identity
coordinates. This assertion is proved locally and tests actual measures.
\end{lemma}
\begin{proof}
Subtract $H(b_*)$ and translate $b_*$. On a small real ball about zero,
the inactive planes have a strict negative gap, so $H(b)=G(b)$ there,
where $G(b)=\max_{p\in P}p\cdot b$ and $P=\partial H(0)$ is the convex
hull of the active slopes. Smooth positive convolution $G_\epsilon$
converges uniformly and has gradients in $P$. If $P$ has nonempty interior,
for any compact $P'\subset\operatorname{Int}P$ there is $a>0$ with
$G(b)-p\cdot b\ge a|b|$ for $p\in P'$. On a fixed ball of radius $r$,
the minimum of $G_\epsilon(b)+\delta|b|^2/2-p\cdot b$ is interior for
sufficiently small $\epsilon,\delta>0$. Its strictly convex gradient is
injective with positive determinant and has a smooth local inverse;
the change-of-variables theorem identifies the determinant integral with
the area of its image. This image contains $P'$ and is contained in
$P+\delta r\overline B$. Let $\delta\downarrow0$ at fixed $\epsilon$,
using continuity of the determinant polynomial and monotonic areas of
outer parallel convex sets, then let $\epsilon\downarrow0$ and exhaust
$\operatorname{Int}P$ by compact inner sets. The determinant integral on
the ball tends to $\operatorname{Area}P$. The boundary of a polygon has
area zero. If $P$ is a segment or point, the same outer-image bound gives
zero in the limit, without an inner-set argument.

Away from zero in this small ball, the local active slopes of the support
function lie on one proper exposed face of $P$, hence on a line. The
local Hessian square vanishes by one-dimensional positive convolution.
Thus the determinant measures converge locally to
$\operatorname{Area}(P)\delta_0$.
On the actual torus chart, the smooth coordinate Hessian square is
supported on its joint Gauss section: off that section a Laurent polynomial
has strict cancellation below its monomial maximum, on an ordinary
neighborhood contained in a tropical hypersurface of dimension at most
one. The coordinate top form is zero there. For the identity moment the
saturated lattice and analytic degree are one;14.5.2--3 and5.10.1(c),(d)
give $2\det\operatorname{Hess}G_\epsilon\,db_1db_2$ on that section.
For every continuous analytic test its restriction to the continuous
Gauss section is continuous. The preceding local weak limit therefore
gives coefficient $2\operatorname{Area}P$ at the actual Gauss point.
Independently, local uniform recursive-product continuity identifies this
limit with the actual Hessian square. Positive Radon reconstruction in
7.2.8--13 supplies all continuous tests. This proves the analytic
calibration, rather than defining MA by a tropical inverse construction.
\end{proof}

\begin{theorem}\label{branch:fanobstruction}
Every $q_H$ belongs to the original $\C$ and satisfies
$\MA(q_H)\ge\delta_g$. Consequently it cannot solve $2\delta_{x_c}$
for any $c>0$. The same obstruction holds for every actual
finite-lower-energy decreasing presentation consisting of such $q_H$
plus constants, and for every weak MA limit of arbitrary families of
these bounded fan weights. It does not exclude solutions elsewhere in
the entire original $\C$ or $\mathcal G$.
\end{theorem}
\begin{proof}
Every active plane slope belongs to $[0,1]^2$: its global supporting
inequality, tested along $b+t e_i$ as $t\to\pm\infty$ with the other
coordinate positive, gives $0\le p_i\le1$. Remove globally redundant
planes. The remaining slopes are nonnegative, so the finite log maximum
is locally a convex maximum of smooth harmonic unit-logarithm slots.
At polynomial zeros the constant lower slot permits omission of all
vanishing terms. Hence the extension is finite continuous and positive
by minimal-sheaf closure. The outside-quadrant equality makes the correction
zero on the same ordinary collar. On the negative quadrant convexity
between points of the two negative axes gives $H\le0=H_0$, so the resulting
weight is normalized and bounded.

The open outside regions have respective affine potentials $b_1$, $b_2$,
and $b_1+b_2$. Their slopes are global supporting slopes for the convex
$H$; since $H(0)=0$, all three lie in $\partial H(0)$. Thus
\[
 \operatorname{conv}\{(1,0),(0,1),(1,1)\}\subset\partial H(0).
\]
This triangle has Euclidean area $1/2$. Near $g$ the translated coordinates
$w_i$ are units; they form an identity coordinate torus, with $g$ its joint
unit Gauss point. Lemma\ref{branch:area} therefore gives
$\MA(q_H)(\{g\})=2\operatorname{Area}(\partial H(0))\ge1$.
This is an actual current coefficient, with the raw $2!$ normalization.
The global positive mass is two, so at most one unit remains away from $g$.
For any weak limit $\lambda$ of these positive measures, nonnegative
continuous tests give $\int h\lambda\ge h(g)$; regularity and Riesz order
give $\lambda\ge\delta_g$. The canonical operator for an already actual
decreasing $\mathcal G$ presentation is its weak MA limit, so it obeys the
same inequality. As $x_c\ne g$, none of these measures equals $2\delta_{x_c}$.
This proves a complete obstruction for all the specified finite fans and
their stated limits, extending the old single-segment obstruction only
within this exact class.
\end{proof}

\subsection{Residual and verification boundary}
The new branch theorem includes arbitrary diffuse and atomic depth
distributions, an exact necessary-and-sufficient full-dual condition,
actual continuous decreasing presentations, a genuine generated operator,
pointwise uniqueness in all actual $\mathcal G$, and canonical nonordered
representation for these targets. It does not change the eligible class
in the original task. The finite-fan obstruction explains why all
outside-quadrant-fixed translated radial modifications retain an atom at
$g$; it is not an obstruction against the entire original class.
General dominated-measure representation, actual general envelopes and
variations, every arbitrary-support equation, and the exact whole-mass
$2\delta_{x_{1/8}}$ equation remain open. Genuine nonalgebraization/descent
and the entire deeper independent audit also remain open. No original
atomic target is claimed solved or refuted. All results are unverified
drafts; finite program checks certify only their declared finite instances.
No compiler or publication is invoked and candidate remains false.
